\chapter{अभिलक्षणिक फलन और केंद्रीय सीमा
प्रमेय}\label{ch:b3:clt}

बृहत् संख्याओं का नियम कहता है कि औसत अभिसरित होते हैं; केंद्रीय
सीमा प्रमेय कहती है कि \emph{वे उतार-चढ़ाव कैसे करते हैं}: $\sqrt n$ से
आवर्धित त्रुटि अनंतस्पर्शी रूप से \omterm{def:b3:clt:gaussianvector}{गाउसीय} होती है --- चाहे आरंभ किसी
भी \omterm{def:b3:probability:space}{बंटन नियम} से किया गया हो। यह सार्वत्रिकता प्रारंभिक प्रायिकता का
सबसे गहरा तथ्य है, और उसकी स्वाभाविक उपपत्ति फूरिये-वैश्लेषिक है:
\emph{\omterm{def:b3:clt:cf}{अभिलक्षणिक फलन}} (किसी \omterm{def:b3:probability:space}{बंटन नियम} का फूरिये रूपांतर) \omterm{def:b3:probability:independence}{स्वतंत्र}
योगों को गुणनफलों में बदल देता है, और \cref{ch:b3:fouriertransform} की मशीनरी --- एकैकीपन,
\omterm{def:b3:clt:gaussianvector}{गाउसीय} स्थिर बिंदु --- इन गुणनफलों के बिंदुशः अभिसरण को बंटन नियमों
के अभिसरण में बदल देती है (लेवी की प्रमेय, पूरी तरह सिद्ध)। अध्याय
\omterm{def:b3:clt:gaussianvector}{गाउसीय} सदिशों पर और सांख्यिकी में सर्वत्र प्रयुक्त विश्वास अंतरालों
की ईमानदार व्युत्पत्ति पर समाप्त होता है; सप्ताहांत समस्या केंद्रीय
सीमा प्रमेय की लिंडेबर्ग वाली दूसरी उपपत्ति स्पष्ट त्रुटि दर के साथ
देती है।

\section{अभिलक्षणिक फलन}

\begin{definition}\label{def:b3:clt:cf}
किसी वास्तविक \omterm{def:b3:probability:space}{यादृच्छिक चर} $X$ का \emph{अभिलक्षणिक
फलन}\index{अभिलक्षणिक फलन}
\[
\varphi_X(\xi) = \E\bigl[\eu^{\iu\xi X}\bigr]
= \int_\R \eu^{\iu\xi x}\,\dd\P_X(x)
\qquad (\xi \in \R)
\]
है (अंतरण प्रमेय उसे \omterm{def:b3:probability:space}{बंटन नियम} से परिकलित कर देती है; किसी \omterm{ex:b3:lebesgue:gamma}{घनत्व}
$f$ के लिए \cref{ch:b3:fouriertransform} की परिपाटी में $\varphi_X(\xi) = \hat f(-\xi)$)।
\end{definition}

\begin{proposition}\label{prop:b3:clt:cfbasics}
(क) $\varphi_X(0) = 1$, $\abs{\varphi_X} \leq 1$, और $\varphi_X$ एकसमान रूप से \omterm{def:b3:topology:continuity}{संतत} है; $\varphi_{aX + b}(\xi) = \eu^{\iu b\xi}\varphi_X(a\xi)$।
(ख) यदि $X, Y$ \emph{\omterm{def:b3:probability:independence}{स्वतंत्र}} हों: $\varphi_{X+Y} = \varphi_X\,\varphi_Y$।
(ग) यदि $\E\abs X^k < \infty$ हो, तो $j \leq
k$ के लिए $\varphi_X^{(j)}(0) = \iu^j\,\E[X^j]$ के साथ $\varphi_X \in \mathcal
C^k$; विशेष रूप से,
प्रसरण $\sigma^2$ वाले केंद्रित $X \in L^2$ के लिए:
\[
\varphi_X(\xi) = 1 - \frac{\sigma^2\xi^2}{2} +
o(\xi^2) \qquad (\xi \to 0).
\]
(घ) \omterm{def:b3:clt:gaussianvector}{गाउसीय}: $X \sim \mathcal N(m, \sigma^2)$ का $\varphi_X(\xi) = \eu^{\iu m\xi - \sigma^2\xi^2/2}$।
\end{proposition}

\begin{proof}
(क) परिबंध तत्काल हैं; और \omterm{def:b3:topology:continuity}{सांतत्य}: प्रभावी अभिसरण से $h
\to 0$ पर $\abs{\varphi(\xi + h)
- \varphi(\xi)} \leq \E\abs{\eu^{\iu hX} - 1} \to 0$,
$\xi$ में एकसमान रूप से। आफ़ीन नियम एक प्रतिस्थापन है। (ख) $\eu^{\iu\xi(X+Y)} =
\eu^{\iu\xi X}\eu^{\iu\xi Y}$, और
\omterm{def:b3:probability:independence}{स्वतंत्र} चरों के गुणनफलों की प्रत्याशाएँ गुणनखंडित हो जाती हैं
(\cref{thm:b3:probability:independence}, वास्तविक और काल्पनिक भागों पर लगाई गई)। (ग) \omterm{def:b3:probability:space}{प्रत्याशा} के
भीतर अवकलन, जिस पर $\E\abs X^j$ का प्रभुत्व है (\cref{thm:b3:lebesgue:paramdiff}); तब $0$ पर टेलर
प्रसार $\mathcal C^2$ फलन $\varphi$ के लिए टेलर--यंग है। (घ) $\mathcal N(0,1)$ के लिए: \omterm{def:b3:clt:gaussianvector}{गाउसीय}
रूपांतर ($a = \frac12$ के साथ \cref{ex:b3:fouriertransform:gaussian}) $\int\eu^{\iu\xi x}\frac{\eu^{-x^2/2}}{\sqrt{2\pi}}\dd
x = \eu^{-\xi^2/2}$ दे देता है; और व्यापक स्थिति आफ़ीन
नियम से।
\end{proof}

\begin{theorem}[एकैकीपन]\label{thm:b3:clt:injectivity}
यदि $\varphi_X = \varphi_Y$ हो, तो $X$ और $Y$ का \omterm{def:b3:probability:space}{बंटन नियम} एक ही है। अधिक
परिशुद्धता से, $X$ से \omterm{def:b3:probability:independence}{स्वतंत्र} $N \sim \mathcal N(0,1)$ और $\varepsilon > 0$ के लिए मृदुकृत चर
$X +
\varepsilon N$ का \omterm{ex:b3:lebesgue:gamma}{घनत्व}
\[
p_\varepsilon(x) = \frac1{2\pi}\int_\R
\varphi_X(-\xi)\,\eu^{-\varepsilon^2\xi^2/2}\,
\eu^{\iu\xi x}\,\dd\xi ,
\]
है, जो केवल $\varphi_X$ से निर्धारित होता है; और $\varepsilon \to 0$ लेने पर $X$ का
\omterm{def:b3:probability:space}{बंटन नियम} पुनः प्राप्त हो जाता है।\index{अभिलक्षणिक फलन, एकैकीपन}
\end{theorem}

\begin{proof}
$X + \varepsilon N$ का \omterm{ex:b3:lebesgue:gamma}{घनत्व} $p_\varepsilon(x) =
\E\bigl[g_\varepsilon(x - X)\bigr]$ है, जहाँ $g_\varepsilon$ $\mathcal N(0, \varepsilon^2)$ \omterm{ex:b3:lebesgue:gamma}{घनत्व} है: वस्तुतः बोरेल
$B$ के लिए \omterm{def:b3:probability:independence}{स्वातंत्र्य} और टोनेली $\P(X + \varepsilon N \in
B) = \int\!\!\int\mathbf 1_B(x + \varepsilon
n)g_1(n)\,\dd n\,\dd\P_X(x) = \int_B\E[g_\varepsilon(t -
X)]\dd t$ दे देते हैं (प्रतिस्थापन
कीजिए, फिर से टोनेली)। $g_\varepsilon$ को उसके रूपांतर के फूरिये प्रतिलोमन से
लिखने पर (\cref{exo:b3:fouriertransform:4}, पुनःमापित): $g_\varepsilon(u) = \frac1{2\pi}\int
\eu^{-\varepsilon^2\xi^2/2}\eu^{\iu\xi u}\dd\xi$, और फुबिनी (सब कुछ \omterm{def:b3:clt:gaussianvector}{गाउसीय} गुणनखंड
से प्रभावित):
\[
p_\varepsilon(x) = \frac1{2\pi}\int_\R
\varphi_X(-\xi)\,\eu^{-\varepsilon^2\xi^2/2}\,
\eu^{\iu\xi x}\,\dd\xi ,
\]
जो केवल $\varphi_X$ का एक फलनक है। यदि $\varphi_X =
\varphi_Y$ हो: तो प्रत्येक $\varepsilon$ के लिए
$X + \varepsilon N$ और $Y + \varepsilon N$ के \omterm{def:b3:probability:space}{बंटन नियम} बराबर हैं; और परिबद्ध \omterm{def:b3:topology:continuity}{संतत} $f$ के लिए
$\varepsilon
\to 0$ पर $\E f(X + \varepsilon N) \to \E f(X)$ (प्रभावी अभिसरण, गुणनफल समष्टि पर बिंदुशः $X + \varepsilon N \to X$), अतः
ऐसे सभी $f$ के लिए $\E f(X) = \E f(Y)$ --- और यह \omterm{def:b3:probability:space}{बंटन नियम} को निर्धारित कर देता
है: प्रत्येक $t$ के लिए $\mathbf 1_{\intoc{-\infty}t}$ को परिबद्ध \omterm{def:b3:topology:continuity}{संतत} ढालों $f_k^\pm$ के बीच
दबाइए (जो $\intoc{-\infty}{t \mp \frac1k}$ पर $1$ के बराबर हैं, $t \pm
\frac1k$ के परे $0$ के, और बीच
में आफ़ीन); $\E
f_k^-(X) \leq F_X(t) \leq \E f_k^+(X)$ में सीमा लेने पर हर उस $t$ पर $F_X(t) =
F_Y(t)$ मिलता है जहाँ
दोनों \omterm{def:b3:topology:continuity}{संतत} हैं, अतः दक्षिण-सांतत्य और उभयनिष्ठ \omterm{def:b3:topology:continuity}{सांतत्य} बिंदुओं की
सघनता से सर्वत्र (दोनों $F$ में गणनीय संख्या के उछाल हैं); और
बराबर बंटन फलन बराबर \omterm{def:b3:probability:space}{बंटन नियम} अनिवार्य कर देते हैं (\cref{exo:b3:measure:3}, जो \cref{thm:b3:measure:uniqueness}
पर टिकी है)।
\end{proof}

\section{बंटन में अभिसरण}

\begin{definition}\label{def:b3:clt:cid}
$X_n$ $X$ की ओर \emph{बंटन में} (अथवा नियम में) अभिसरित होता है,
जिसे $X_n \Rightarrow X$ लिखा जाता है, यदि\index{बंटन में अभिसरण}
\[
\E\bigl[f(X_n)\bigr] \longrightarrow \E\bigl[f(X)\bigr]
\qquad\text{प्रत्येक परिबद्ध संतत } f\colon\R\to\R \text{ के लिए} .
\]
तुल्य रूप से (\cref{exo:b3:clt:4}): $F_X$ के प्रत्येक \omterm{def:b3:topology:continuity}{सांतत्य} बिंदु $t$ पर $F_{X_n}(t) \to F_X(t)$।
$X_n$ का किसी उभयनिष्ठ \omterm{def:b3:probability:space}{प्रायिकता समष्टि} पर रहना आवश्यक नहीं: केवल
\omterm{def:b3:probability:space}{बंटन नियम} मायने रखते हैं।
\end{definition}

\begin{theorem}[हेली की चयन प्रमेय]
\label{thm:b3:clt:helly}
बंटन फलनों के प्रत्येक अनुक्रम $(F_n)$ का कोई ऐसा उपानुक्रम होता है
जो सीमा के प्रत्येक \omterm{def:b3:topology:continuity}{सांतत्य} बिंदु पर किसी अनह्रासमान दक्षिण-संतत
$G \colon
\R \to \intcc01$ की ओर बिंदुशः अभिसरित होता है --- संभवतः $G(+\infty) - G(-\infty) <
1$ के साथ (द्रव्यमान
अनंत की ओर पलायन कर सकता है)।\index{हेली की प्रमेय}
\end{theorem}

\begin{proof}
विकर्ण निष्कर्षण प्रत्येक परिमेय $q$ के लिए $F_{n_k}(q) \to \ell(q)$ दे देता है
(मान \omterm{def:b3:topology:compact}{संहत} $\intcc01$ में हैं)। $G(t) = \inf\{\ell(q) : q \in \Q, q > t\}$ परिभाषित कीजिए: यह अनह्रासमान है;
और दक्षिण-संतत भी (दाईं ओर से सिकुड़ते परिमेय प्रतिवेशों पर एक
निम्नतम)। $G$ के किसी \omterm{def:b3:topology:continuity}{सांतत्य} बिंदु $t$ पर: परिमेय $q_1 < t < q_2$ के
लिए प्रत्येक $F_{n_k}$ की एकदिष्टता से
\[
\ell(q_1) \leq \liminf F_{n_k}(t) \leq \limsup F_{n_k}(t)
\leq \ell(q_2),
\]
और $G$ की निम्नतम वाली परिभाषा तथा परिमेयों पर $\ell$ की एकदिष्टता
से: जब भी $s < q$ हो तब $G(s) \leq \ell(q) \leq G(q)$। $s < q_1 < t$ लेने पर $\ell(q_1) \geq G(s)$ और $\ell(q_2) \leq G(q_2)$ मिलते
हैं; और $s \uparrow t$ तथा $q_2
\downarrow t$ लेने पर $t$ पर $G$ की \omterm{def:b3:topology:continuity}{सांतत्य} $\liminf$ और
$\limsup$ दोनों को $G(t)$ तक दबा देती है।
\end{proof}

\begin{lemma}[अभिलक्षणिक फलन से सुसंहतता]
\label{lem:b3:clt:tightness}
किसी भी \omterm{def:b3:probability:space}{यादृच्छिक चर} $X$ और $u > 0$ के लिए:
\[
\P\Bigl(\abs X \geq \frac2u\Bigr) \;\leq\;
\frac1u\int_{-u}^{u}\bigl(1 -
\operatorname{Re}\varphi_X(\xi)\bigr)\,\dd\xi .
\]
\end{lemma}

\begin{proof}
टोनेली--फुबिनी से (समाकल्य परिबद्ध, क्षेत्र $\xi$ में परिमित):
\[
\frac1u\int_{-u}^u\bigl(1 -
\operatorname{Re}\varphi_X(\xi)\bigr)\dd\xi
= \E\Bigl[\frac1u\int_{-u}^u(1 - \cos(\xi X))\,\dd\xi\Bigr]
= 2\,\E\Bigl[1 - \frac{\sin(uX)}{uX}\Bigr]
\]
(कोष्ठक को $X = 0$ पर उसकी सीमा $0$ मानिए)। समाकल्य ऋणेतर है
($\abs{\sin t} \leq \abs t$), और $\abs{uX} \geq 2$ के लिए: $1 - \frac{\sin(uX)}{uX} \geq 1 -
\frac1{\abs{uX}} \geq \frac12$। अतः \omterm{def:b3:probability:space}{प्रत्याशा} के भीतर केवल घटना
$\{\abs{uX} \geq 2\}$ रखने पर कम से कम $2 \cdot \frac12\,\P(\abs X \geq \frac2u)$ बच जाता है, और यही दावा है।
\end{proof}

\begin{theorem}[लेवी की सांतत्य प्रमेय]
\label{thm:b3:clt:levy}
मान लीजिए $(X_n)$ ऐसे \omterm{def:b3:probability:space}{यादृच्छिक चर} हैं जिनके \omterm{def:b3:clt:cf}{अभिलक्षणिक फलन} बिंदुशः
अभिसरित होते हैं: प्रत्येक $\xi$ के लिए $\varphi_{X_n}(\xi) \to
\varphi(\xi)$, जहाँ $\varphi =
\varphi_X$ किसी
\omterm{def:b3:probability:space}{यादृच्छिक चर} $X$ का \omterm{def:b3:clt:cf}{अभिलक्षणिक फलन} है। तब $X_n \Rightarrow X$।\index{लेवी की
सांतत्य प्रमेय}
\end{theorem}

\begin{proof}
\emph{सुसंहतता।} $\varepsilon > 0$ स्थिर कीजिए। चूँकि $\varphi$ $\varphi(0) = 1$ के साथ $0$
पर \omterm{def:b3:topology:continuity}{संतत} है, अतः $\frac1u\int_{-u}^u(1 - \operatorname{Re}\varphi) <
\varepsilon$ वाला $u > 0$ चुनिए; प्रभावी अभिसरण से (स्थिर
$[-u,u]$ पर समाकल्य $2$ से परिबद्ध) बड़े $n$ के लिए $\varphi_{X_n}$ का वही
समाकल $< 2\varepsilon$ है: \cref{lem:b3:clt:tightness} बड़े $n$ के लिए $\P(\abs{X_n} \geq \frac2u)
\leq 2\varepsilon$ दे देता है, और अचर
को बड़ा करने पर शेष परिमित संख्या के मामले भी सँभल जाते हैं: अर्थात्
\omterm{def:b3:probability:space}{बंटन नियम} \emph{सुसंहत} हैं --- कोई द्रव्यमान पलायन नहीं करता।

\emph{उपानुक्रम।} मान लीजिए $(F_{n_k})$ कोई उपानुक्रम है; हेली (\cref{thm:b3:clt:helly})
से \omterm{def:b3:topology:continuity}{सांतत्य} बिंदुओं पर $F_{n_{k_j}} \to G$ निकालिए। सुसंहतता $G(-\infty) = 0$, $G(+\infty) = 1$
अनिवार्य कर देती है (\omterm{def:b3:topology:continuity}{सांतत्य} बिंदुओं पर $G(\frac2u) - G(-\frac2u) \geq 1 -
2\varepsilon$): अतः $G$ किसी
\omterm{def:b3:probability:space}{यादृच्छिक चर} $Y$ का सच्चा बंटन फलन है। तब $X_{n_{k_j}} \Rightarrow Y$ (\cref{exo:b3:clt:4}, $F$
से बंटन-संबंधी अभिसरण), अतः \emph{बिंदुशः} $\varphi_{X_{n_{k_j}}} \to \varphi_Y$ ($x
\mapsto \eu^{\iu\xi x}$ परिबद्ध
\omterm{def:b3:topology:continuity}{संतत} है, वास्तविक और काल्पनिक भाग पृथक-पृथक); और परिकल्पना से तुलना
करने पर $\varphi_Y = \varphi = \varphi_X$, तथा एकैकीपन (\cref{thm:b3:clt:injectivity}) $Y \sim X$ दे देता है, अर्थात् $G =
F_X$।

\emph{निष्कर्ष।} $(F_n)$ के प्रत्येक उपानुक्रम का कोई उप-उपानुक्रम
\emph{उसी} $F_X$ की ओर (उसके \omterm{def:b3:topology:continuity}{सांतत्य} बिंदुओं पर) अभिसरित होता है;
अतः प्रत्येक \omterm{def:b3:topology:continuity}{सांतत्य} बिंदु $t$ पर $F_n(t) \to F_X(t)$ (वह वास्तविक अनुक्रम
अभिसरित होता है जिसके सभी उपानुक्रमों के उप-उपानुक्रमों की सीमा एक
ही हो): अर्थात् $X_n \Rightarrow X$।
\end{proof}

\section{केंद्रीय सीमा प्रमेय}

\begin{theorem}[केंद्रीय सीमा प्रमेय]\label{thm:b3:clt:clt}
मान लीजिए $(X_n)$ $\E X_1 = m$ और $\V(X_1) =
\sigma^2 \in \intoo0\infty$ वाले \omterm{def:b3:probability:independence}{स्वतंत्र} सम-बंटित चर हैं।
तब\index{केंद्रीय सीमा प्रमेय}
\[
\frac{S_n - nm}{\sigma\sqrt n} \;\Longrightarrow\; \mathcal
N(0, 1) :
\qquad
\P\Bigl(a \leq \frac{S_n - nm}{\sigma\sqrt n} \leq
b\Bigr) \longrightarrow
\frac{1}{\sqrt{2\pi}}\int_a^b\eu^{-x^2/2}\,\dd x
\]
सभी $a < b$ के लिए।
\end{theorem}

\begin{proof}
केंद्रीकरण और मानकीकरण कीजिए: $Z_i = \frac{X_i - m}{\sigma}$ (\omterm{def:b3:probability:independence}{स्वतंत्र} सम-बंटित, माध्य $0$,
प्रसरण $1$) और $T_n =
\frac1{\sqrt n}\sum_{i\leq n}Z_i$। \omterm{def:b3:probability:independence}{स्वातंत्र्य} तथा आफ़ीन नियम (\cref{prop:b3:clt:cfbasics}) से:
\[
\varphi_{T_n}(\xi) =
\varphi_{Z}\Bigl(\frac{\xi}{\sqrt n}\Bigr)^{n},
\qquad
\varphi_Z(\eta) = 1 - \frac{\eta^2}2 + \eta^2\rho(\eta),\quad
\rho(\eta)\to0 .
\]
$\xi$ स्थिर कीजिए और $a_n = \varphi_Z(\xi/\sqrt n)$, $b_n = 1 -
\frac{\xi^2}{2n}$ लीजिए: बड़े $n$ के लिए दोनों
का मापांक $\leq 1$ है ($\xi^2 \leq 4n$ होते ही $\abs{b_n} \leq 1$; और $\abs{a_n} \leq 1$ सदैव)। $\abs a, \abs b \leq 1$ के
लिए प्रारंभिक असमिका $\abs{a^n - b^n} \leq
n\abs{a - b}$ (दूरबीनी $a^n - b^n = \sum a^k(a - b)b^{n-1-k}$)
\[
\Bigl|\varphi_{T_n}(\xi) - \Bigl(1 -
\frac{\xi^2}{2n}\Bigr)^{n}\Bigr|
\leq n\,\Bigl|\varphi_Z\Bigl(\frac\xi{\sqrt n}\Bigr) - 1 +
\frac{\xi^2}{2n}\Bigr|
= \xi^2\,\Bigl|\rho\Bigl(\frac{\xi}{\sqrt n}\Bigr)\Bigr|
\longrightarrow 0,
\]
दे देती है, जबकि $\bigl(1 - \frac{\xi^2}{2n}\bigr)^n \to
\eu^{-\xi^2/2}$ (वास्तविक लघुगणक)। अतः प्रत्येक $\xi$ के लिए
$\varphi_{T_n}(\xi) \to
\eu^{-\xi^2/2} = \varphi_{\mathcal N(0,1)}(\xi)$ (\cref{prop:b3:clt:cfbasics}(घ)): और लेवी (\cref{thm:b3:clt:levy}) से निष्कर्ष निकलता है $T_n \Rightarrow \mathcal
N(0,1)$।
अंतराल प्रायिकताएँ इसलिए निकलती हैं क्योंकि $F_{\mathcal
N}$ सर्वत्र \omterm{def:b3:topology:continuity}{संतत} है।
\end{proof}

\begin{example}[विश्वास अंतराल, ईमानदारी से व्युत्पन्न]
\label{ex:b3:clt:confidence}
$n$ \omterm{def:b3:probability:independence}{स्वतंत्र} मतदाताओं का सर्वेक्षण कीजिए; $\hat p_n = S_n/n$ $\sigma^2 = p(1-p) \leq \frac14$ के साथ
वास्तविक $p$ का आकलन करता है। केंद्रीय सीमा प्रमेय बड़े $n$ के
लिए
\[
\P\Bigl(\abs{\hat p_n - p} \leq
\frac{z}{2\sqrt n}\Bigr)
\;\geq\; \P\Bigl(\Bigl|\frac{S_n - np}{\sigma\sqrt n}\Bigr|
\leq z\Bigr)
\longrightarrow \Phi(z) - \Phi(-z),
\]
देती है, जहाँ $\Phi$ मानक \omterm{def:b3:clt:gaussianvector}{गाउसीय} बंटन फलन है। $z = 1.96$ के साथ:
अनंतस्पर्शी विश्वास $95\%$, और सीमांत $\frac{1.96}{2\sqrt n} \leq 3\%$ के लिए $n \geq
\bigl(\frac{1.96}{0.06}\bigr)^2 \approx 1068$ चाहिए ---
यही वह संख्या है जो पढ़ने में आने वाले हर ``$\pm3$ अंक, $95\%$'' के
पीछे है; इसकी तुलना चेबिशेव वाले $5556$ (\cref{exo:b3:probability:7}) से कीजिए। $\sqrt n$
सार्वत्रिक है: त्रुटि आधी करने के लिए प्रतिदर्श चौगुना कीजिए ---
यही वह नियम है जो \omterm{ex:b3:probability:sllnapps}{मोंते कार्लो} की लागत भी तय करता है (\cref{exo:b3:clt:7})।
\end{example}

\section{गाउसीय सदिश}

\begin{definition}\label{def:b3:clt:gaussianvector}
कोई यादृच्छिक सदिश $X = (X_1, \dots, X_d)$ \emph{गाउसीय}\index{गाउसीय सदिश} कहलाता
है यदि प्रत्येक रैखिक संचय $\langle t, X\rangle = \sum t_iX_i$ कोई (संभवतः अपह्रासी) वास्तविक
गाउसीय चर हो। उसका \omterm{def:b3:probability:space}{बंटन नियम} माध्य सदिश $m = (\E X_i)$ और \emph{सहप्रसरण
आव्यूह} $\Sigma = \bigl(\operatorname{Cov}
(X_i, X_j)\bigr)$ से निर्धारित होता है: वस्तुतः सदिश का \omterm{def:b3:clt:cf}{अभिलक्षणिक फलन}
$\varphi_X(t) = \E\eu^{\iu\langle t, X\rangle}$ $\langle t, X\rangle$ के \omterm{def:b3:clt:cf}{अभिलक्षणिक फलन} का $1$ पर मान है:
\[
\varphi_X(t) = \exp\Bigl(\iu\langle t, m\rangle -
\tfrac12\,t^{\mathsf T}\Sigma\,t\Bigr),
\]
और $d$-विमीय \omterm{def:b3:clt:cf}{अभिलक्षणिक फलन} एकैकी होते हैं (वही मृदुकरण वाली
उपपत्ति जो \cref{thm:b3:clt:injectivity} में है, निर्देशांकशः गाउसीय के साथ)।
\end{definition}

\begin{theorem}\label{thm:b3:clt:gaussianvector}
मान लीजिए $X$ कोई \omterm{def:b3:clt:gaussianvector}{गाउसीय सदिश} है।
\begin{enumerate}
  \item प्रत्येक आफ़ीन प्रतिबिंब $AX + b$ \omterm{def:b3:clt:gaussianvector}{गाउसीय सदिश} है।
  \item घटक $X_i$ \emph{\omterm{def:b3:probability:independence}{स्वतंत्र}} हैं तभी और केवल तभी जब $\Sigma$
        विकर्ण हो: संयुक्त रूप से \omterm{def:b3:clt:gaussianvector}{गाउसीय} चरों के लिए, असहसंबंधित
        $=$ \omterm{def:b3:probability:independence}{स्वतंत्र}।
  \item यदि $\Sigma$ व्युत्क्रमणीय हो, तो $X$ का \omterm{ex:b3:lebesgue:gamma}{घनत्व} $\frac{1}{(2\pi)^{d/2}\sqrt{\det\Sigma}}
        \exp\bigl(-\frac12(x - m)^{\mathsf T}\Sigma^{-1}(x -
        m)\bigr)$ है।
\end{enumerate}
\end{theorem}

\begin{proof}
(1) $AX + b$ के घटकों के रैखिक संचय $X$ के रैखिक संचयों के आफ़ीन फलन
हैं: अतः \omterm{def:b3:clt:gaussianvector}{गाउसीय} (किसी \omterm{def:b3:clt:gaussianvector}{गाउसीय} चर का आफ़ीन प्रतिबिंब \omterm{def:b3:clt:gaussianvector}{गाउसीय} होता है)।
(2) यदि $\Sigma$ विकर्ण हो, तो \omterm{def:b3:clt:cf}{अभिलक्षणिक फलन} गुणनखंडित हो जाता है:
$\varphi_X(t) = \prod_i\exp(\iu t_im_i -
\frac12\Sigma_{ii}t_i^2) = \prod\varphi_{X_i}(t_i)$, जो गुणनफल \omterm{def:b3:probability:space}{बंटन नियम} का \omterm{def:b3:clt:cf}{अभिलक्षणिक फलन} है (\cref{thm:b3:probability:independence} को
$d$-विमीय एकैकीपन के माध्यम से पढ़िए): अतः घटक \omterm{def:b3:probability:independence}{स्वतंत्र} हैं।
विलोम \omterm{def:b3:probability:independence}{स्वतंत्र} $L^2$ चरों के सहप्रसरणों का लुप्त होना है।
(3) $\Sigma = P D P^{\mathsf T}$ का विकर्णीकरण कीजिए ($P$ लांबिक, $D > 0$ विकर्ण ---
\cref{exo:b3:submanifolds:8}); सदिश $Y = P^{\mathsf T}(X - m)$ सहप्रसरण $D$ वाला \omterm{def:b3:clt:gaussianvector}{गाउसीय} है: (2) से उसके घटक
\omterm{def:b3:probability:independence}{स्वतंत्र} $\mathcal N(0, d_i)$ हैं, अतः $Y$ का \omterm{ex:b3:lebesgue:gamma}{घनत्व} गुणनफल \omterm{ex:b3:lebesgue:gamma}{घनत्व} है; इसे
आयतन-संरक्षी $x = m + PY$ (\cref{thm:b3:product:linearchange}, $\abs{\det P} = 1$) से अग्रसारित कीजिए और चरघातांक
को निश्चर रूप में फिर से लिखिए।
\end{proof}

\begin{theorem}[बहुविमीय केंद्रीय सीमा प्रमेय]
\label{thm:b3:clt:multiclt}
मान लीजिए $(X_n)$ माध्य $m$ और सहप्रसरण आव्यूह $\Sigma$ वाले $\R^d$ के
\omterm{def:b3:probability:independence}{स्वतंत्र} सम-बंटित वर्ग-समाकलनीय यादृच्छिक \emph{सदिश} हैं। तब $\frac{S_n - nm}{\sqrt n}$
\omterm{def:b3:clt:gaussianvector}{गाउसीय सदिश} $\mathcal N(0,
\Sigma)$ की ओर बंटन में अभिसरित होता
है।\index{केंद्रीय सीमा प्रमेय!बहुविमीय}
\admitted
\end{theorem}

\begin{remark}\label{rem:b3:clt:multiclt}
लगभग सब कुछ पहले से ही हमारे हाथ में है। प्रत्येक दिशा $t \in \R^d$ के लिए
वास्तविक चर $\langle t,
\frac{S_n - nm}{\sqrt n}\rangle$ प्रसरण $t^{\mathsf T}\Sigma t$ वाले \omterm{def:b3:probability:independence}{स्वतंत्र} सम-बंटित वास्तविक चरों
का मानकीकृत योग है, अतः \cref{thm:b3:clt:clt} का परिकलन $d$-विमीय अभिलक्षणिक फलनों
का $\eu^{-t^{\mathsf T}\Sigma t/2}$ की ओर बिंदुशः अभिसरण दे देता है, जो $\mathcal N(0, \Sigma)$ (\cref{def:b3:clt:gaussianvector}) का
\omterm{def:b3:clt:cf}{अभिलक्षणिक फलन} है। जो हमने पुनः सिद्ध नहीं किया वह है लेवी की
\omterm{def:b3:topology:continuity}{सांतत्य} प्रमेय \emph{$\R^d$ में}: हेली का चयन और सुसंहतता आकलन
सामान्य रीति से (निर्देशांकशः) सामान्यीकृत हो जाते हैं, और यह
क्रामेर--वोल्ड न्यूनीकरण किसी भी स्नातकोत्तर प्रायिकता पाठ्यक्रम
में ईमानदारी से किया जाता है; इस अध्याय की विधियों से बाहर कुछ भी
आवश्यक नहीं।
\end{remark}

\begin{method}\label{met:b3:clt:method}
किसी सीमा \omterm{def:b3:probability:space}{बंटन नियम} की पहचान के लिए: \omterm{def:b3:clt:cf}{अभिलक्षणिक फलन} परिकलित कीजिए,
बिंदुशः सीमा लीजिए, उसे पहचानिए (\omterm{def:b3:clt:gaussianvector}{गाउसीय} $\eu^{-\sigma^2\xi^2/2}$, प्वासों $\eu^{\lambda(\eu^{\iu\xi}-1)}$,
चरघातांकी $\frac{\lambda}
{\lambda - \iu\xi}$, \dots) और लेवी लगाइए। यह तीन-चरणीय अनुष्ठान
(\omterm{def:b3:probability:independence}{स्वातंत्र्य} $\to$ गुणनफल; $0$ पर टेलर $\to$ चरघातांकी सीमा; लेवी
$\to$ नियम में अभिसरण) केंद्रीय सीमा प्रमेय, विरल घटनाओं का प्वासों
नियम (\cref{exo:b3:clt:5}), और इस पाठ्यक्रम की प्रत्येक शास्त्रीय सीमा प्रमेय
सिद्ध कर देता है। लगभग निश्चित कथनों के लिए \cref{ch:b3:probability} के साधन-संदूक पर
लौटिए: दोनों अध्याय एक ही $S_n$ के बारे में भिन्न प्रश्नों का उत्तर
देते हैं।
\end{method}

\section{अभ्यास}

\begin{exercise}[$\star$]\label{exo:b3:clt:1}
\omterm{def:b3:clt:cf}{अभिलक्षणिक फलन} परिकलित कीजिए: $\intcc{-1}1$ पर एकसमान; चरघातांकी $\mathcal E(\lambda)$;
प्वासों $\mathcal P(\lambda)$; द्विपद $\mathcal B(n, p)$। \cref{thm:b3:clt:injectivity} के माध्यम से निकालिए कि \omterm{def:b3:probability:independence}{स्वतंत्र}
प्वासों चरों ($\lambda, \mu$) का योग प्वासों $(\lambda +
\mu)$ है।
\end{exercise}

\begin{exercise}[$\star\star$]\label{exo:b3:clt:2}
(क) दर्शाइए कि $\varphi_X$ वास्तविक-मान वाला है तभी और केवल तभी जब $X$
और $-X$ का \omterm{def:b3:probability:space}{बंटन नियम} एक ही हो (कोई \emph{सममित} चर)।
(ख) मान लीजिए किसी $\xi_0 \neq
0$ के लिए $\abs{\varphi_X(\xi_0)} = 1$। दर्शाइए कि $X$ लगभग निश्चित
रूप से किसी समांतर श्रेढ़ी $a + \frac{2\pi}{\xi_0}\Z$ पर आलंबित है \emph{($\varphi_X(\xi_0) = \eu^{\iu\theta}$ लिखिए
और $\E[1 -
\cos(\xi_0X - \theta)]$ परिकलित कीजिए)}। निकालिए कि यदि $X$ का कोई \omterm{ex:b3:lebesgue:gamma}{घनत्व} हो, तो
सभी $\xi \neq 0$ के लिए $\abs{\varphi_X(\xi)} < 1$।
\end{exercise}

\begin{exercise}[$\star\star$]\label{exo:b3:clt:3}
मान लीजिए $X \sim \mathcal N(m_1, \sigma_1^2)$ और $Y \sim
\mathcal N(m_2, \sigma_2^2)$ \omterm{def:b3:probability:independence}{स्वतंत्र} हैं। दर्शाइए कि $X + Y
\sim \mathcal N(m_1 + m_2, \sigma_1^2 + \sigma_2^2)$, और अधिक
व्यापक रूप से यह कि \omterm{def:b3:clt:gaussianvector}{गाउसीय} कुल \omterm{def:b3:probability:independence}{स्वतंत्र} योगों तथा आफ़ीन प्रतिचित्रणों
के अंतर्गत स्थायी है। विरोधाभास: क्या दो \emph{आश्रित} \omterm{def:b3:clt:gaussianvector}{गाउसीय} चरों
का योग सदैव \omterm{def:b3:clt:gaussianvector}{गाउसीय} होता है? (\cref{exo:b3:clt:9}।)
\end{exercise}

\begin{exercise}[$\star\star$]\label{exo:b3:clt:4}
(क) \cref{def:b3:clt:cid} में दी गई तुल्यता सिद्ध कीजिए: यदि सभी परिबद्ध \omterm{def:b3:topology:continuity}{संतत} $f$
के लिए $\E f(X_n) \to \E f(X)$ हो, तो \omterm{def:b3:topology:continuity}{सांतत्य} बिंदुओं पर $F_{X_n}(t) \to F_X(t)$ \emph{($\mathbf 1_{\intoc{-\infty}t}$ को दो
\omterm{def:b3:topology:continuity}{संतत} सीढ़ी-ढालों के बीच दबाइए)}; और विलोमतः \emph{(किसी परिबद्ध
\omterm{def:b3:topology:continuity}{संतत} $f$ का ढाल फलनों के योगों से आसन्नन कीजिए, अथवा \omterm{def:b3:topology:continuity}{सांतत्य}
बिंदुओं के किसी सूक्ष्म जालक पर प्रतिबंधन कीजिए)} --- विलोम को पहले
एकसमान \omterm{def:b3:topology:continuity}{संतत} $f$ के लिए और फिर व्यापक रूप से लिया जा सकता है।
(ख) दर्शाइए कि $X_n \Rightarrow c$ (कोई अचर) से प्रायिकता में $X_n
\to c$ निकलता है।
\end{exercise}

\begin{exercise}[$\star\star$]\label{exo:b3:clt:5}
(विरल घटनाओं का नियम) मान लीजिए $np_n \to \lambda > 0$ के साथ $X_n \sim \mathcal B(n, p_n)$। अभिलक्षणिक
फलनों और \cref{thm:b3:clt:levy} के माध्यम से दर्शाइए कि $X_n \Rightarrow \mathcal
P(\lambda)$। संख्यात्मक विवेक
जाँच: $\mathcal B(100, 0.02)$ और $\mathcal P(2)$ के लिए $\P(X = 0)$ की तुलना कीजिए।
\end{exercise}

\begin{exercise}[$\star\star$]\label{exo:b3:clt:6}
(क) कोई निष्पक्ष पासा $n = 1000$ बार फेंका जाता है; कुल के $3600$ से अधिक
होने की प्रायिकता का आसन्नन कीजिए (माध्य $3500$, प्रति फेंक प्रसरण
$\frac{35}{12}$)।
(ख) $S \sim \mathcal B(100, \frac12)$ के लिए \omterm{def:b3:topology:continuity}{सांतत्य} संशोधन ($\pm\frac12$) के साथ केंद्रीय सीमा प्रमेय
से $\P(45 \leq S \leq 55)$ का आसन्नन कीजिए, और संशोधन के प्रभाव पर टिप्पणी कीजिए।
\end{exercise}

\begin{exercise}[$\star\star$]\label{exo:b3:clt:7}
(\omterm{ex:b3:probability:sllnapps}{मोंते कार्लो} त्रुटि) \cref{pb:b3:probability:1}, प्रश्न 11 की व्यवस्था में, $g \in
L^2(\intcc01^d)$ के
साथ मान लीजिए $\sigma^2 = \V(g(U_1))$ और $I = \int
g$। दर्शाइए
\[
\sqrt n\,\Bigl(\frac1n\sum_{k\leq n}g(U_k) - I\Bigr)
\Longrightarrow \mathcal N(0, \sigma^2),
\]
और अनंतस्पर्शी $95\%$ त्रुटि-दंड $\pm
1.96\,\sigma/\sqrt n$ निकालिए --- जो विमा $d$ से
\omterm{def:b3:probability:independence}{स्वतंत्र} है। इसकी तुलना विमा $d$ में नियतांकी मध्यबिंदु नियम से
कीजिए ($\mathcal C^2$ समाकल्यों के लिए त्रुटि $\sim n^{-2/d}$): किस विमा से आगे
यादृच्छिक प्रतिचयन जीतने लगता है?
\end{exercise}

\begin{exercise}[$\star\star\star$]\label{exo:b3:clt:8}
(स्लुत्स्की) मान लीजिए $X_n \Rightarrow X$ और प्रायिकता में $Y_n \to c$ ($c$ अचर)।
दर्शाइए कि $X_n + Y_n \Rightarrow X +
c$ और $Y_nX_n \Rightarrow cX$। \emph{(अभिलक्षणिक फलनों और परिबंध $\abs{\E\eu^{\iu\xi
(X_n+Y_n)} - \eu^{\iu\xi c}\E\eu^{\iu\xi X_n}} \leq
\E\abs{\eu^{\iu\xi(Y_n - c)} - 1}$
के साथ काम कीजिए, $\abs{Y_n - c}
\leq \delta$ पर विभाजन कीजिए।)} अनुप्रयोग: \cref{ex:b3:clt:confidence} में
अज्ञात $\sigma = \sqrt{p(1-p)}$ को $\sqrt{\hat p_n(1 - \hat
p_n)}$ से बदलना उचित ठहराइए।
\end{exercise}

\begin{exercise}[$\star\star\star$]\label{exo:b3:clt:9}
मान लीजिए $X \sim \mathcal N(0,1)$ और $\P(\varepsilon = \pm1) = \frac12$ वाला $\varepsilon$ \omterm{def:b3:probability:independence}{स्वतंत्र} हैं; $Y =
\varepsilon X$ रखिए।
(क) दर्शाइए कि $Y \sim \mathcal N(0,1)$ और $\operatorname{Cov}(X, Y) = 0$।
(ख) दर्शाइए कि $X$ और $Y$ \omterm{def:b3:probability:independence}{स्वतंत्र} \emph{नहीं} हैं, और यह कि
$(X, Y)$ \omterm{def:b3:clt:gaussianvector}{गाउसीय सदिश} नहीं है \emph{($\P(X + Y =
0)$ परिकलित कीजिए)}।
(ग) शिक्षा: \cref{thm:b3:clt:gaussianvector}(2) के लिए संयुक्त गाउसीयता आवश्यक है ---
``असहसंबंधित \omterm{def:b3:clt:gaussianvector}{गाउसीय}'' अकेले कुछ भी सिद्ध नहीं करता।
\end{exercise}

\begin{exercise}[$\star\star$]\label{exo:b3:clt:10}
कोशी \omterm{def:b3:probability:space}{बंटन नियम} का \omterm{ex:b3:lebesgue:gamma}{घनत्व} $\frac1{\pi(1 + x^2)}$ है।
(क) दर्शाइए कि उसका \omterm{def:b3:clt:cf}{अभिलक्षणिक फलन} $\eu^{-\abs\xi}$ है (\cref{exo:b3:fouriertransform:1} और प्रतिलोमन)।
(ख) दर्शाइए कि यदि $X_1, \dots, X_n$ \omterm{def:b3:probability:independence}{स्वतंत्र} सम-बंटित कोशी हों, तो $\frac{S_n}n$ पुनः
कोशी है --- और \emph{वही} \omterm{def:b3:probability:space}{बंटन नियम}: औसत कभी संकेंद्रित नहीं होता।
(ग) इसका बृहत् संख्याओं के नियमों तथा केंद्रीय सीमा प्रमेय से मेल
बिठाइए: कौन-सी परिकल्पनाएँ विफल होती हैं? ($\E\abs{X_1}$ परिकलित कीजिए।)
\end{exercise}

\begin{exercise}[$\star\star$]\label{exo:b3:clt:11}
(बीज रूप में स्थायी \omterm{def:b3:probability:space}{बंटन नियम}) मान लीजिए $(X_n)$ \omterm{def:b3:probability:independence}{स्वतंत्र} सम-बंटित
मानक कोशी हैं (\cref{exo:b3:clt:10})।
(क) दर्शाइए कि किसी भी $a, b > 0$ के लिए $aX_1 + bX_2$ का \omterm{def:b3:probability:space}{बंटन नियम} $(a + b)X_1$ का
है: अर्थात् कोशी कुल सूचकांक $1$ का \emph{कड़ाई से स्थायी} है।
(ख) दर्शाइए कि \omterm{def:b3:clt:gaussianvector}{गाउसीय} कुल सूचकांक $2$ का कड़ाई से स्थायी है:
\omterm{def:b3:probability:independence}{स्वतंत्र} सम-बंटित $\mathcal N(0,1)$ चरों $X_i$ के लिए $aX_1 + bX_2 \sim \sqrt{a^2 + b^2}\,X_1$।
(ग) $\eu^{-c\abs\xi^\alpha}$ रूप के अभिलक्षणिक फलनों के माध्यम से समझाइए कि सूचकांक-$\alpha$
स्थायित्व योगों के लिए मानकीकरण $n^{1/\alpha}$ क्यों अनिवार्य कर देता है,
और यह केंद्रीय सीमा प्रमेय के आकर्षण-द्रोणियों के बारे में क्या
कहता है: कौन-से \omterm{def:b3:probability:independence}{स्वतंत्र} सम-बंटित योग आफ़ीन मानकीकरण के बाद \omterm{def:b3:clt:gaussianvector}{गाउसीय}
के बजाय किसी कोशी \omterm{def:b3:probability:space}{बंटन नियम} की ओर अभिसरित हो सकते हैं?
\end{exercise}

\begin{exercise}[$\star\star$]\label{exo:b3:clt:12}
(आनुभविक बंटन फलन) मान लीजिए $(X_n)$ बंटन फलन $F$ वाले \omterm{def:b3:probability:independence}{स्वतंत्र}
सम-बंटित चर हैं, और $F_n(t) =
\frac1n\#\{k \leq n : X_k \leq t\}$।
(क) $t$ स्थिर कीजिए। दर्शाइए कि $n F_n(t) \sim \mathcal B(n, F(t))$, कि लगभग निश्चित रूप से
$F_n(t) \to F(t)$ (\cref{thm:b3:probability:slln}), और यह कि
\[
\sqrt n\,\bigl(F_n(t) - F(t)\bigr) \Longrightarrow
\mathcal N\bigl(0,\ F(t)(1 - F(t))\bigr) .
\]
(ख) किस $t$ पर अनंतस्पर्शी प्रसरण उच्चिष्ठ है? व्याख्या कीजिए:
माध्यिका वही स्थान है जहाँ किसी आनुभविक बंटन को ठीक-ठीक बाँधना सबसे
कठिन है।
(ग) \omterm{def:b3:topology:continuity}{संतत} $F$ के लिए दर्शाइए कि $\sup_t\abs{F_n(t) - F(t)}$ का \omterm{def:b3:probability:space}{बंटन नियम} $F$ पर निर्भर
नहीं करता \emph{(\cref{exo:b3:probability:1} के माध्यम से एकसमान चरों तक न्यूनीकरण
कीजिए)} --- यही वह बंटन-निरपेक्ष चमत्कार है जो कोल्मोगोरोव--स्मिरनोव
परीक्षण के पीछे है; उस \omterm{def:b3:probability:space}{बंटन नियम} का कोई परिकलन नहीं माँगा जा रहा।
\end{exercise}

\section{समस्या: केंद्रीय सीमा प्रमेय की लिंडेबर्ग वाली उपपत्ति,
दर के साथ}

\begin{problem}[{सप्ताहांत समस्या --- प्रतिस्थापन विधि}]
\label{pb:b3:clt:1}
लिंडेबर्ग (1922) ने केंद्रीय सीमा प्रमेय एक निहत्था कर देने वाली
सरल कल्पना से सिद्ध की: \emph{योज्य पदों को एक-एक करके \omterm{def:b3:clt:gaussianvector}{गाउसीय} चरों
से बदल दीजिए} और प्रत्येक अदला-बदली को टेलर प्रसार से नियंत्रित
कीजिए। इस विधि को फूरिये विश्लेषण की आवश्यकता नहीं, यह स्पष्ट त्रुटि
दर देती है, और आज समूचे प्रायिकता सिद्धांत में सार्वत्रिकता की
उपपत्तियों को चलाती है। मान लीजिए $(X_i)$ \omterm{def:b3:probability:independence}{स्वतंत्र} सम-बंटित, केंद्रित,
$\V(X_1) = 1$ वाले, $\beta = \E\abs{X_1}^3 < \infty$ के साथ हैं; और मान लीजिए $(N_i)$ $X_i$ से \omterm{def:b3:probability:independence}{स्वतंत्र}
\omterm{def:b3:probability:independence}{स्वतंत्र} सम-बंटित $\mathcal N(0,1)$ हैं (अस्तित्व: \cref{thm:b3:probability:existence})।
\[
T_n = \frac{X_1 + \dots + X_n}{\sqrt n},
\qquad
G_n = \frac{N_1 + \dots + N_n}{\sqrt n} \sim \mathcal N(0,1).
\] रखिए।

\medskip
\noindent\textbf{भाग I --- अदला-बदली सर्वसमिका।} $f
\in \mathcal C^3_b(\R)$ स्थिर कीजिए
(तीन परिबद्ध \omterm{def:b3:topology:continuity}{संतत} अवकलज; $M_3 = \sup\abs{f'''}$)। $0 \leq i \leq n$ के लिए संकर योग
\[
H_i = \frac{X_1 + \dots + X_i + N_{i+1} + \dots +
N_n}{\sqrt n},
\]
परिभाषित कीजिए, जिससे $H_n = T_n$ और $H_0 = G_n$।
\begin{enumerate}
  \item $W_i =
        \frac{1}{\sqrt n}\bigl(\sum_{j<i}X_j +
        \sum_{j>i}N_j\bigr)$ के साथ $H_i = W_i + \frac{X_i}{\sqrt n}$ और $H_{i-1}
        = W_i + \frac{N_i}{\sqrt n}$ लिखिए, और ध्यान दीजिए कि $W_i$
        युग्म $(X_i, N_i)$ से \omterm{def:b3:probability:independence}{स्वतंत्र} है। इसे उचित ठहराइए।
  \item समाकल या लाग्रांज शेषफल के साथ टेलर: किन्हीं भी वास्तविक
        $w, h$ के लिए:
        \[
        \Bigl|f(w + h) - f(w) - f'(w)h -
        \tfrac12f''(w)h^2\Bigr| \leq
        \frac{M_3\,\abs h^3}{6} .
        \]
  \item प्रश्न 2 को दो बार लगाइए ($w = W_i$ पर $h = \frac{X_i}{\sqrt n}$ और $h = \frac{N_i}{\sqrt n}$),
        \omterm{def:b3:probability:space}{प्रत्याशा} लीजिए, और \omterm{def:b3:probability:independence}{स्वातंत्र्य} तथा $X_i$ और $N_i$ के पहले
        दो आघूर्णों के मिलान का उपयोग करके दर्शाइए
        \[
        \bigl|\E f(H_i) - \E f(H_{i-1})\bigr|
        \leq \frac{M_3}{6}\cdot
        \frac{\beta + \gamma}{n^{3/2}},
        \qquad \gamma = \E\abs{N_1}^3 =
        \frac{2\sqrt2}{\sqrt\pi} .
        \]
  \item $i$ पर दूरबीनी योग लीजिए और \emph{लिंडेबर्ग परिबंध}
        निष्कर्ष निकालिए:
        \[
        \bigl|\E f(T_n) - \E f(G_n)\bigr| \leq
        \frac{M_3\,(\beta + \gamma)}{6\,\sqrt n} .
        \]
\end{enumerate}

\medskip
\noindent\textbf{भाग II --- चिकने $f$ से केंद्रीय सीमा प्रमेय
तक।}
\begin{enumerate}[resume]
  \item दर्शाइए कि प्रत्येक $f \in
        \mathcal C_b^3$ के लिए $\E f(T_n) \to \E f(N)$, और इसे सभी परिबद्ध
        \omterm{def:b3:topology:continuity}{संतत} $f$ तक उन्नत कीजिए: ऐसा $f$ और $\varepsilon$ दिए जाने पर
        किसी बड़े अंतराल पर
        $\norm{f - f_\varepsilon}_\infty \leq \varepsilon$
        वाला $f_\varepsilon \in \mathcal C^3_b$ रचिए --- उदाहरणार्थ $f$ का किसी $\mathcal C^\infty$ उभार
        (\cref{thm:b3:lp:regularization}) के साथ संवलन कीजिए --- और पुच्छों को सुसंहतता
        ($\V(T_n) = 1$ और चेबिशेव) से सँभालिए। निष्कर्ष निकालिए $T_n \Rightarrow \mathcal N(0, 1)$:
        अर्थात् केंद्रीय सीमा प्रमेय, पुनः सिद्ध।
  \item उपपत्ति ने कहाँ इसका उपयोग किया कि $X_i$ \emph{सम}-बंटित
        हैं? दर्शाइए कि उसने मुश्किल से ही किया: $\sum_i\V(X_i) = s_n^2$ और तृतीय
        आघूर्णों वाले \omterm{def:b3:probability:independence}{स्वतंत्र}, केंद्रित, असम $X_i$ के लिए वह रूप
        कहिए और सिद्ध कीजिए, और त्रुटि $\frac{M_3}{6s_n^3}\sum_i\bigl(\E\abs{X_i}^3
        + \V(X_i)^{3/2}\gamma\bigr)$ प्राप्त कीजिए ---
        अर्थात् लिंडेबर्ग की सच्ची प्रमेय अपने ल्यापुनोव रूप में।
\end{enumerate}

\medskip
\noindent\textbf{भाग III --- मात्रात्मक लाभांश।}
\begin{enumerate}[resume]
  \item (बंटन फलन) मान लीजिए $t \in \R$ और $\mathbf 1_{\intoc{-\infty}t}$ का ऊपर तथा नीचे से
        चौड़ाई $\delta$ के $\mathcal C^3_b$ ढालों से आसन्नन कीजिए (उन्हें रचिए,
        $M_3 = O(\delta^{-3})$ के साथ)। भाग I के साथ मिलाकर स्पष्ट अचरों वाला
        द्विपद परिबंध
        \[
        \sup_{t\in\R}\,\bigl|\P(T_n \leq t) -
        \Phi(t)\bigr| \;\leq\;
        \frac{C_1(\beta + \gamma)}{\delta^3\sqrt n} +
        C_2\,\delta
        \qquad (\text{प्रत्येक } \delta > 0),
        \]
        निकालिए ($C_2\delta$ पद इस बात का उपयोग करता है कि $\Phi$ का \omterm{ex:b3:lebesgue:gamma}{घनत्व}
        $\frac1{\sqrt{2\pi}}$ से परिबद्ध है), और $n^{-1/8}$ कोटि की एकसमान दर प्राप्त
        करने के लिए $\delta \sim
        n^{-1/8}$ पर अनुकूलन कीजिए।
        (अनुकूलतम $n^{-1/2}$ --- बेरी--एसेन --- के लिए सूक्ष्मतर साधन
        चाहिए; यहाँ बात यह है कि प्रारंभिक अदला-बदली से भी एक
        \emph{स्पष्ट} दर मिल जाती है।)
  \item (दे म्वाव्र--लाप्लास, मात्रात्मक रूप में) $X_i = 2B_i - 1$ (निष्पक्ष
        सिक्कों के चिह्न) पर विशेषीकरण कीजिए: निष्कर्ष की तुलना
        \cref{pb:b3:product:1}, प्रश्न 7 के स्थानीय आकलन से कीजिए --- प्रत्येक विधि
        वह क्या देती है जो दूसरी नहीं देती?
  \item (सार्वत्रिकता) एक अनुच्छेद में समझाइए कि प्रतिस्थापन विधि
        केंद्रीय सीमा प्रमेय से अधिक क्यों दिखाती है: चिकने $f$
        के साथ $\E f(\text{योग})$ रूप की कोई भी सांख्यिकी, $n^{-1/2}$ की कोटि पर,
        योज्य पदों के पहले दो आघूर्णों से परे उनके \emph{समूचे
        \omterm{def:b3:probability:space}{बंटन नियम}} के प्रति असंवेदी है --- यही वह ``निश्चरता
        सिद्धांत'' है जो आधुनिक सार्वत्रिकता परिणामों (यादृच्छिक
        आव्यूह, यादृच्छिक बहुपद) के मूल में है और जिसका पहला
        उदाहरण केंद्रीय सीमा प्रमेय है।
\end{enumerate}

\medskip
\noindent\textbf{भाग IV --- मृदुकरण, और आगे बढ़ाकर: बेहतर दरें।}
$n^{-1/2}$ (चिकने $f$) से $n^{-1/8}$ (बंटन फलन) तक की हानि $f'''$ को
उच्चतम मानक में लगाने से आई थी। संकर उसका कुछ भाग सुधार सकते हैं:
उनमें \omterm{def:b3:clt:gaussianvector}{गाउसीय} योज्य पद होते हैं, और \omterm{def:b3:clt:gaussianvector}{गाउसीय} \emph{मृदु करते} हैं।
\begin{enumerate}[resume]
  \item (एक छिपा हुआ \omterm{def:b3:clt:gaussianvector}{गाउसीय}) $1 \leq i \leq n - 1$, $h =
        \frac{X_i}{\sqrt n}$ अथवा $\frac{N_i}{\sqrt n}$, और $\theta \in \intcc01$ के
        लिए $Z = \frac{N_{i+1} + \dots + N_n}{\sqrt
        n}$ के साथ $W_i + \theta h = A +
        Z$ लिखिए। दर्शाइए कि $Z \sim \mathcal N\bigl(0,
        \frac{n-i}n\bigr)$ युग्म $(A,
        h)$
        से \omterm{def:b3:probability:independence}{स्वतंत्र} है, और प्रत्येक \omterm{def:b3:topology:continuity}{संतत} $g \in
        L^1(\R)$ के लिए निकालिए
        \[
        \E\bigl[\abs h^3\,\abs{g(W_i + \theta h)}\bigr]
        \;\leq\; \sqrt{\frac{n}{2\pi(n - i)}}\;
        \norm{g}_{L^1}\;\E\abs h^3 .
        \]
  \item प्रश्न 10 को टेलर शेषफल के समाकल रूप
        \[
        f(w + h) = f(w) + f'(w)h + \tfrac12f''(w)h^2 +
        \int_0^1\frac{(1 - \theta)^2}2\,f'''(w + \theta
        h)\,h^3\,\dd\theta,
        \]
        के साथ मिलाकर प्रश्न 3--4 फिर से कीजिए: $f \in \mathcal C^3_b$ के लिए, और
        साथ में $f''' \in L^1(\R)$ होने पर,
        \[
        \bigl|\E f(T_n) - \E f(G_n)\bigr| \leq
        \frac{\beta + \gamma}{3\sqrt{2\pi}}\cdot
        \frac{\norm{f'''}_{L^1}}{\sqrt n} +
        \frac{M_3(\beta + \gamma)}{6\,n^{3/2}}
        \]
        \emph{(प्रश्न 10 अदला-बदलियों $i \leq n - 1$ को सँभालता है ---
        $\sum_{m=1}^{n-1}m^{-1/2} \leq 2\sqrt n$ का उपयोग कीजिए --- और अंतिम को प्रश्न 3 का कच्चा
        परिबंध सँभालता है)}। जाँचिए कि प्रश्न 7 के ढाल
        $\norm{\psi_\delta'''}_{L^1} = K_1\delta^{-2}$
        संतुष्ट करते हैं जबकि $M_3 = K\delta^{-3}$, उन्हें भीतर डालिए, और $\delta$
        पर अनुकूलन कीजिए: एकसमान बंटन-फलन दर सुधरकर $O(n^{-1/6})$ हो जाती
        है।
  \item (एक और आघूर्ण का मिलान) अतिरिक्त रूप से $\E
        X_1^3 = 0$ और $\beta_4 = \E X_1^4 < \infty$
        मान लीजिए। $\E N_1^3$ और $\E N_1^4$ परिकलित कीजिए, चतुर्थ कोटि तक
        प्रसार कीजिए, और उन्हीं पंक्तियों पर सिद्ध कीजिए कि
        बंटन-फलन दर $O(n^{-1/4})$ हो जाती है \emph{(अब $\norm{\psi_\delta^{(4)}}_{L^1} =
        K_2\delta^{-3}$ और $M_4 = K'\delta^{-4}$;
        $\delta = n^{-1/4}$ चुनिए)}।
  \item (बाधा) मान लीजिए $X_1$ के पहले $k$ आघूर्ण \omterm{def:b3:clt:gaussianvector}{गाउसीय}
        आघूर्णों से मेल खाते हैं ($k = 2$ सदैव; $k = 3$ ठीक तभी जब
        $\E X_1^3 = 0$; और $k \geq 4$ वस्तुतः कभी नहीं, क्योंकि $\E N_1^4 = 3$)। सत्यापित
        कीजिए कि $\delta^{-k}n^{-(k-1)/2}$ को $\delta$ के विरुद्ध संतुलित करके प्रश्न
        10--12 की योजना बंटन-फलन दर $n^{-(k-1)/(2k+2)}$ दे देती है, और ध्यान
        दीजिए कि यह घातांक बेरी--एसेन मान $\frac12$ तक केवल $k \to
        \infty$ पर
        ही पहुँचता है। कुछ वाक्यों में समझाइए कि अदला-बदली विधि
        क्यों संतृप्त हो जाती है: प्रत्येक अदला-बदली निरपेक्ष मान
        में लगाई जाती है, जबकि फूरिये मार्ग (एसेन की मृदुकरण
        असमिका) अभिलक्षणिक-फलन अंतर के दोलन का दोहन करता है और
        केवल तीन आघूर्णों से $C\beta n^{-1/2}$ तक पहुँच जाता है।
\end{enumerate}

\medskip
\noindent\textbf{भाग V --- दो विमाएँ: बहुविमीय केंद्रीय सीमा
प्रमेय, अदला-बदली से।} अब मान लीजिए $X_i$ सहप्रसरण आव्यूह $\Sigma$ और
$\beta' = \E\norm{X_1}^3 <
\infty$ (यूक्लिडीय मानक) वाले $\R^2$ के \omterm{def:b3:probability:independence}{स्वतंत्र} सम-बंटित केंद्रित
यादृच्छिक \emph{सदिश} हैं।
\begin{enumerate}[resume]
  \item (\omterm{def:b3:clt:gaussianvector}{गाउसीय सदिश}, कोटि तक) $\Sigma =
        PDP^{\mathsf T}$ का विकर्णीकरण कीजिए (\cref{exo:b3:submanifolds:8})
        और $C = P\sqrt DP^{\mathsf T}$ रखिए। \omterm{def:b3:probability:independence}{स्वतंत्र} मानक \omterm{def:b3:clt:gaussianvector}{गाउसीय} चरों के किसी युग्म
        $Z = (Z^1,
        Z^2)$ (\cref{thm:b3:probability:existence}) के लिए दर्शाइए कि $N
        = CZ$ माध्य $0$, सहप्रसरण
        $\Sigma$ और $\gamma' = \E\norm N^3 <
        \infty$ वाला \omterm{def:b3:clt:gaussianvector}{गाउसीय सदिश} (\cref{def:b3:clt:gaussianvector}) है; और यह कि
        \omterm{def:b3:probability:independence}{स्वतंत्र} सम-बंटित प्रतिकृतियों $N_i$ के लिए $G_n = \frac{N_1 + \dots +
        N_n}{\sqrt n}$ का \omterm{def:b3:probability:space}{बंटन
        नियम} \emph{ठीक-ठीक} $\mathcal N(0, \Sigma)$ है।
  \item (दो चरों में टेलर) $M_3 =
        \max_{\abs\alpha = 3}\sup\abs{\partial^\alpha f} <
        \infty$ वाले $\mathcal C^3$ वर्ग के $f \colon \R^2 \to
        \R$ के लिए
        सिद्ध कीजिए
        \[
        \Bigl|f(w + h) - f(w) - \langle\nabla f(w),
        h\rangle - \tfrac12\langle h, D^2f(w)\,h\rangle
        \Bigr| \leq \frac{M_3}6\,\bigl(\abs{h_1} +
        \abs{h_2}\bigr)^3 \leq \frac{\sqrt2\,M_3}3\,
        \norm h^3
        \]
        \emph{($\intcc01$ पर $t \mapsto f(w + th)$ का अध्ययन कीजिए)}।
  \item ($\R^2$ में केंद्रीय सीमा प्रमेय) सदिश संकरों $H_i$ पर
        प्रतिस्थापन योजना चलाइए: दर्शाइए कि प्रथम और द्वितीय कोटि
        के पद निरस्त हो जाते हैं (माध्य और सहप्रसरण मेल खाते हैं),
        दूरबीनी योग लीजिए, और प्रश्न 5 की भाँति उन्नत कीजिए ($\E\norm{T_n}^2 =
        \operatorname{tr}\Sigma$
        से सुसंहतता; और मृदुकरण अब $\R^2$, \cref{thm:b3:lp:regularization} में) ताकि निष्कर्ष
        निकले: प्रत्येक परिबद्ध \omterm{def:b3:topology:continuity}{संतत} $f \colon
        \R^2 \to \R$ के लिए
        \[
        \E\,f\Bigl(\frac{X_1 + \dots + X_n}{\sqrt n}\Bigr)
        \longrightarrow \E\,f(N), \qquad N \sim \mathcal
        N(0, \Sigma) :
        \]
        अर्थात् विमा $2$ में \cref{thm:b3:clt:multiclt}, चिकने $f$ के लिए एक दर के
        साथ और बिना किसी फूरिये विश्लेषण के।
  \item (क्रामेर--वोल्ड, और एक संयुक्त उतार-चढ़ाव) निकालिए कि
        प्रत्येक स्थिर $t \in \R^2$ के लिए $\langle t, \frac{S_n}{\sqrt n}\rangle
        \Rightarrow \mathcal N(0, t^{\mathsf T}\Sigma t)$। अनुप्रयोग: \omterm{def:b3:probability:independence}{स्वतंत्र}
        सम-बंटित वास्तविक $(\xi_i)$ के लिए, जो केंद्रित हैं, $\E\xi_1^2 = 1$,
        $\E\xi_1^6 < \infty$ (ताकि भाग V $V_i = (\xi_i, \xi_i^2 - 1)$ पर लागू हो), दर्शाइए
        \[
        \frac1{\sqrt n}\Bigl(\sum_{i\leq n}\xi_i,\
        \sum_{i\leq n}(\xi_i^2 - 1)\Bigr) \Longrightarrow
        \mathcal N\Bigl(0, \begin{pmatrix} 1 & \E\xi_1^3\\
        \E\xi_1^3 & \E\xi_1^4 - 1\end{pmatrix}\Bigr) :
        \]
        अर्थात् आनुभविक माध्य और आनुभविक द्वितीय आघूर्ण संयुक्त
        रूप से \omterm{def:b3:clt:gaussianvector}{गाउसीय} ढंग से उतार-चढ़ाव करते हैं --- और सीमा में
        \omterm{def:b3:probability:independence}{स्वतंत्र} रूप से तभी और केवल तभी जब $\E\xi_1^3 = 0$ (\cref{thm:b3:clt:gaussianvector})।
\end{enumerate}

\medskip
\noindent\textbf{भाग VI --- डेल्टा विधि।}
\begin{enumerate}[resume]
  \item मान लीजिए $(\hat\theta_n)$ ऐसे \omterm{def:b3:probability:space}{यादृच्छिक चर} हैं कि किसी वास्तविक
        प्राचल $\theta$ के लिए $\sqrt n(\hat\theta_n - \theta) \Rightarrow
        \mathcal N(0, \sigma^2)$, और मान लीजिए $g$ $\theta$ पर
        अवकलनीय है। \emph{डेल्टा विधि} सिद्ध कीजिए:
        \[
        \sqrt n\bigl(g(\hat\theta_n) - g(\theta)\bigr)
        \Longrightarrow \mathcal N\bigl(0,
        g'(\theta)^2\sigma^2\bigr)
        \]
        \emph{($\theta$ पर $\eta \to 0$ के साथ $g(x) - g(\theta) = (g'(\theta) +
        \eta(x))(x - \theta)$ लिखिए; प्रायिकता में
        $\hat\theta_n \to \theta$ और फिर $\eta(\hat\theta_n) \to 0$ दर्शाइए; और स्लुत्स्की, \cref{exo:b3:clt:8} तथा
        \cref{exo:b3:clt:4}(ख) से समाप्त कीजिए)}।
  \item अनुप्रयोग। (क) माध्य $\mu$ और प्रसरण $\sigma^2$ वाले \omterm{def:b3:probability:independence}{स्वतंत्र}
        सम-बंटित वास्तविक $(\xi_i)$ और $\bar X_n =
        \frac1n\sum_{i\leq n}\xi_i$ के लिए: दर्शाइए कि $\mu \neq 0$
        होने पर $\sqrt n(\bar
        X_n^2 - \mu^2) \Rightarrow \mathcal N(0,
        4\mu^2\sigma^2)$, और यह कि $\mu = 0$ के लिए सही कथन किसी अन्य
        मापक्रम पर रहता है: $N \sim \mathcal N(0,1)$ के साथ $n\bar X_n^2 \Rightarrow \sigma^2N^2$ (सीमा का बंटन फलन
        पहचानिए)। (ख) (प्रसरण स्थिरीकरण) किसी $\mathcal B(1, p)$ प्रतिदर्श की
        सफलता आवृत्ति $\hat p_n$ और $p \in \intoo01$ के लिए दर्शाइए कि $g(p) = \arcsin\sqrt p$
        $p$ \emph{कुछ भी हो}
        \[
        \sqrt n\,\bigl(g(\hat p_n) - g(p)\bigr)
        \Longrightarrow \mathcal N\Bigl(0, \frac14\Bigr)
        \]
        संतुष्ट करता है --- अर्थात् अज्ञात प्राचल से मुक्त एक
        अनंतस्पर्शी त्रुटि-दंड; इसकी तुलना \cref{ex:b3:clt:confidence} से कीजिए।
\end{enumerate}

\medskip
\noindent\textbf{भाग VII --- प्वासों, उसी विधि से: ले काम की
प्रमेय।} प्रतिस्थापन एक दूसरा सार्वत्रिकता वर्ग भी जानता है: बहुत
सारी \omterm{def:b3:probability:independence}{स्वतंत्र} \emph{विरल} घटनाओं के योग। $\N$ पर बंटन नियमों के
लिए उपयुक्त दूरी \emph{कुल विचरण} है,
\[
d_{\mathrm{TV}}(\mu, \nu) = \sup_{A\subseteq\N}\,
\abs{\mu(A) - \nu(A)} .
\]
\begin{enumerate}[resume]
  \item दर्शाइए कि $d_{\mathrm{TV}}(\mu, \nu) =
        \frac12\sum_{k\geq0}\abs{\mu(\{k\}) -
        \nu(\{k\})}$, और युग्मन परिबंध सिद्ध कीजिए: एक ही
        समष्टि पर $\mu$ और $\nu$ बंटन नियमों वाले यादृच्छिक चरों
        के \emph{किसी भी} युग्म $(X, Y)$ के लिए $d_{\mathrm{TV}}(\mu, \nu) \leq \P(X \neq Y)$।
  \item $p \in \intoo01$ के लिए ठीक-ठीक परिकलित कीजिए:
        \[
        d_{\mathrm{TV}}\bigl(\mathcal B(1, p), \mathcal
        P(p)\bigr) = p\bigl(1 - \eu^{-p}\bigr) \leq p^2 .
        \]
  \item (ले काम, अदला-बदली से) मान लीजिए $X_i \sim \mathcal B(1,
        p_i)$ और $Y_i \sim \mathcal P(p_i)$, जिनमें
        $2n$ चर \omterm{def:b3:probability:independence}{स्वतंत्र} हैं; $S = X_1 + \dots + X_n$, और स्मरण कीजिए कि $\lambda = \sum_ip_i$
        (\cref{exo:b3:clt:1}) के साथ $Y_1 + \dots + Y_n \sim \mathcal P(\lambda)$। पूर्णांक संकरों $H_i = Y_1 + \dots + Y_i + X_{i+1} + \dots
        + X_n$ में एक बार में
        एक निर्देशांक बदलिए: प्रत्येक $A \subseteq \N$ के लिए दर्शाइए
        \[
        \abs{\P(H_{i-1} \in A) - \P(H_i \in A)} \leq
        d_{\mathrm{TV}}\bigl(\mathcal B(1, p_i), \mathcal
        P(p_i)\bigr),
        \]
        और \emph{ले काम असमिका} निष्कर्ष निकालिए:
        \[
        d_{\mathrm{TV}}\bigl(S \text{ का बंटन नियम},\ \mathcal
        P(\lambda)\bigr) \leq \sum_{i=1}^np_i^2 .
        \]
  \item लाभांश। (क) $p_i = \frac\lambda n$ के लिए: परिबंध $\frac{\lambda^2}n$ है --- अर्थात् विरल
        घटनाओं का नियम (\cref{exo:b3:clt:5}) एक स्पष्ट दर तक उन्नत, जो सभी
        घटनाओं पर एकसमान है और असमान $p_i$ के लिए भी मान्य। (ख)
        $500$ पत्र वितरित किए जाते हैं, और प्रत्येक \omterm{def:b3:probability:independence}{स्वतंत्र} रूप
        से प्रायिकता $\frac1{500}$ के साथ भटक जाता है: प्राचल $1$ वाले
        प्वासों प्रतिरूप की त्रुटि परिबद्ध कीजिए, और किसी भी पत्र
        के न भटकने की प्रायिकता का आकलन कीजिए। (ग) समस्या समाप्त
        कीजिए: यहाँ मिले दोनों सार्वत्रिकता वर्गों की तुलना कीजिए
        --- \omterm{def:b3:clt:gaussianvector}{गाउसीय} (बहुत सारे छोटे बिखरे हुए योगदान; दो आघूर्णों
        का मिलान; टेलर) और प्वासों (बहुत सारे विरल योगदान; एक
        माध्य का मिलान; एक यथार्थ कुल-विचरण युग्मन) --- और दोनों
        के पीछे खड़ी एक ही प्रतिस्थापन विधि की भी।

  \item (सापेक्ष त्रुटि और लघुगणक रूपांतर) मान लीजिए $(X_n)$
        \omterm{def:b3:probability:independence}{स्वतंत्र} सम-बंटित, धनात्मक, माध्य $\mu > 0$, प्रसरण $\sigma^2$ वाले
        हैं, और $\bar X_n$ आनुभविक माध्य। डेल्टा विधि से दर्शाइए कि
        \[
        \sqrt n\,\bigl(\ln\bar X_n - \ln\mu\bigr)
        \Longrightarrow
        \mathcal N\Bigl(0,\ \frac{\sigma^2}{\mu^2}\Bigr) :
        \]
        अर्थात् $\ln\bar X_n$ का अनंतस्पर्शी प्राचल \emph{विचरण गुणांक}
        $\sigma/\mu$ है --- सापेक्ष, मापन-मुक्त त्रुटि। $\mu$ के लिए गुणक
        रूप $\bar X_n\cdot\eu^{\pm1.96\,\sigma/(\mu\sqrt
        n)}$ वाला $95\%$ विश्वास अंतराल निकालिए, और समझाइए कि
        वह योगात्मक अंतराल से कब बेहतर है।
  \item (तृतीय आघूर्ण त्रुटि को दिशा देता है) केंद्रित बर्नूली($p$)
        $X \sim$ के लिए $\E\bigl[(X -
        p)^3\bigr] = p(1-p)(1-2p)$ परिकलित कीजिए। भाग IV के विश्लेषण (अदला-बदली
        त्रुटि तृतीय आघूर्णों से चालित होती है) का उपयोग करते हुए
        समझाइए कि $p \neq \frac12$ के लिए $\mathcal B(n, p)$ का प्रसामान्य आसन्नन असममित
        क्यों होता है --- एक ओर अधिक और दूसरी ओर कम --- और $p = \frac12$
        तीव्रतर मिलित-आघूर्ण दर का आनंद क्यों लेता है। $\mathcal N(2, 1.8)$ के
        विरुद्ध $\mathcal
        B(20, 0.1)$ पर विषमता का चिह्न संख्यात्मक रूप से सत्यापित
        कीजिए: $\P(S = 0) = 0.9^{20}$ की तुलना $\intoo{-\infty}{0.5}$ के \omterm{def:b3:clt:gaussianvector}{गाउसीय} द्रव्यमान से कीजिए।
\end{enumerate}
\end{problem}
